\section*{Capítulo \ref{ch:b1:complexation} --- Complejación}

\begin{solution}{exo:b1:complexation:1}
$[\ce{Cu(NH3)4}]^{2+}$: \ce{Cu^2+}, cuatro \ce{NH3} unidos por N.
$[\ce{Fe(CN)6}]^{4-}$: \ce{Fe^2+}, seis \ce{CN-} unidos por C;
$[\ce{Fe(CN)6}]^{3-}$: lo mismo con \ce{Fe^3+}. $[\ce{Ag(NH3)2}]^{+}$:
\ce{Ag+}, dos \ce{NH3} unidos por N. $[\ce{FeSCN}]^{2+}$: \ce{Fe^3+}, un
\ce{SCN-}, unido por N a pesar de cómo se escribe la fórmula.
$[\ce{Zn(OH)4}]^{2-}$: \ce{Zn^2+}, cuatro \ce{OH-} unidos por O.
$[\ce{CaY}]^{2-}$: \ce{Ca^2+}, un ion EDTA unido por dos N y cuatro O.
\end{solution}

\begin{solution}{exo:b1:complexation:2}
$\log K_i = \log\beta_i - \log\beta_{i-1}$: 4.10, 3.41, 2.82, 2.03, y
$\mathrm{p}K_{d,i} = \log K_i$ toma los mismos valores. El último se
refiere a
$[\ce{Cu(NH3)4}]^{2+} \rightleftharpoons [\ce{Cu(NH3)3}]^{2+} + \ce{NH3}$.
\end{solution}

\begin{solution}{exo:b1:complexation:3}
Dominios: $[\ce{Cu(NH3)4}]^{2+}$ por debajo de pNH3 2.03; después, $n = 3$
hasta 2.82; $n = 2$ hasta 3.41; $n = 1$ hasta 4.10, y \ce{Cu^2+} por
encima. pNH3 = 0:
$[\ce{Cu(NH3)4}]^{2+}$; 2.5: $[\ce{Cu(NH3)3}]^{2+}$; 3.5:
$[\ce{CuNH3}]^{2+}$; 5.0: \ce{Cu^2+}.
\end{solution}

\begin{solution}{exo:b1:complexation:4}
$[\ce{FeSCN^2+}]/[\ce{Fe^3+}] = \beta_1[\ce{SCN-}] = 10^{2.96} \times 0.10 =
91$: la fracción complejada es $91/92 = \qty{98.9}{\percent}$.
\end{solution}

\begin{solution}{exo:b1:complexation:5}
Cada $[\mathrm{ML}_k] = \beta_k[\mathrm{M}][\mathrm{L}]^k$; al dividir
por la suma sobre $k$ (el metal total) se elimina $[\mathrm{M}]$. Con tres
especies, $f_i = 1/(1 + [\mathrm{ML}_{i-1}]/[\mathrm{ML}_i] +
[\mathrm{ML}_{i+1}]/[\mathrm{ML}_i]) = 1/(1 + 1/(K_i x) + K_{i+1}x)$ con
$x = [\mathrm{L}]$. El denominador es mínimo cuando se anula su derivada,
$-1/(K_i x^2) + K_{i+1}$, es decir, $x^2 = 1/(K_iK_{i+1})$, donde los dos
vecinos son iguales, y $\mathrm{pL} = \frac12(\log K_i + \log K_{i+1})$.
Para $[\ce{Cu(NH3)2}]^{2+}$: $\frac12(3.41 + 2.82) = 3.12$.
\end{solution}

\begin{solution}{exo:b1:complexation:6}
$[\ce{FeSCN}]^{2+} + \ce{F-} \rightleftharpoons [\ce{FeF}]^{2+} + \ce{SCN-}$,
$K = 10^{6.85 - 2.96} = 10^{3.89}$. Cociente $= K[\ce{F-}]/[\ce{SCN-}] =
10^{3.89} \times 0.010/0.10 = 780$: el \omterm{def:b1:complexation:complex}{complejo} rojo queda sustituido
casi por completo y el color desaparece.
\end{solution}

\begin{solution}{exo:b1:complexation:7}
$[\ce{MgY}]^{2-} + \ce{Ca^2+} \rightleftharpoons [\ce{CaY}]^{2-} + \ce{Mg^2+}$,
$K = 10^{12.69 - 10.90} = 10^{1.79} = 62$. Con cantidades iguales,
$x^2/(1 - x)^2 = 62$, $x/(1 - x) = 7.9$, $x = 0.89$: se libera el
\qty{89}{\percent} del magnesio.
\end{solution}

\begin{solution}{exo:b1:complexation:8}
\ce{Ag2O(s) + H2O + 4NH3 <=> 2[Ag(NH3)2]+ + 2OH-}, $K = \beta_2^2 \times
10^{-15.43} = 10^{14.44 - 15.43} = 10^{-0.99}$. La constante no es
pequeña: en cuanto el amoniaco libre alcanza unas décimas de mol por
litro, el óxido pardo formado por las primeras gotas (en las que el
amoniaco actúa como base) se disuelve dando el \omterm{def:b1:complexation:complex}{complejo} amoniacal de la
plata.
\end{solution}

\begin{solution}{exo:b1:complexation:9}
pH 10: $\alpha = 1 + 10^{1.24} + 10^{-1.96} = 18.4$, $\log\beta' = 12.69 -
1.26 = 11.43$. pH 7: $\alpha = 1 + 10^{4.24} + 10^{4.04} + 10^{0.19} +
\cdots = 10^{4.45}$, $\log\beta' = 8.24$. A pH 7 la constante condicional
es inferior a $10^{10}$: la reacción del calcio con el EDTA no es lo
bastante completa para una valoración, porque la mayor parte del EDTA está
protonada.
\end{solution}

\begin{solution}{exo:b1:complexation:10}
1.~\ce{CuO(s) + H2O + 4NH3 <=> [Cu(NH3)4]^2+ + 2OH-}, $K = \beta_4 \times
10^{-20.64} = 10^{12.36 - 20.64} = 10^{-8.28}$.
2.~$\mathrm{pH} = 7 + \frac12(9.25 + \log 1.0) = 11.63$, $[\ce{OH-}] =
10^{-2.37}$, y $[\ce{Cu(NH3)4^2+}] = 10^{-8.28}/10^{-4.75} =
\qty{3.0e-4}{mol/L}$: se disuelve muy poco.
3.~$\mathrm{pH} = 9.25$, $[\ce{OH-}] = 10^{-4.75}$, y la fórmula da
$10^{-8.28 + 9.50} = 10^{1.22}$, mucho más de lo que el amoniaco puede
transportar: el óxido se disuelve hasta que se agota el amoniaco, no
hasta el equilibrio. Los iones amonio consumen el hidróxido que libera la
disolución; por eso las sales de amonio ayudan al amoniaco a disolver los
compuestos de cobre.
\end{solution}

\begin{solution}{exo:b1:complexation:11}
1.~La reacción es la segunda formación menos la primera: $K = K_2/K_1 =
10^{3.90 - 3.32} = 10^{0.58} = 3.8 > 1$, de modo que $[\ce{AgNH3}]^{+}$
reacciona consigo mismo cuando es la especie principal.
2.~Por el \cref{exo:b1:complexation:5}, la fracción máxima está en
$\mathrm{pNH_3} = \frac12(3.32 + 3.90) = 3.61$, donde $\beta_1 x = 10^{-0.29}
= 0.51$ y $\beta_2x^2 = 1.0$: $f = 0.51/(1 + 0.51 + 1.0) = 0.20$. Nunca
más del \qty{20}{\percent} de la plata está como $[\ce{AgNH3}]^{+}$.
\end{solution}

\begin{solution}{exo:b1:complexation:12}
1.~\ce{CaCO3(s) + Y^4- <=> [CaY]^2- + CO3^2-}, $K = \beta K_s = 10^{12.69 -
8.30} = 10^{4.39}$.
2.~$x^2/(0.10 - x) = 10^{4.39}$ da $x = \qty{0.10}{mol/L}$ con un margen
de $4 \times 10^{-7}$: se gasta todo el EDTA, \qty{10.0}{g} de carbonato
de calcio por litro.
3.~En medio ácido, el EDTA se protona ($\mathrm{p}K_{a4} = 11.24$) y su
constante condicional baja; el carbonato también se protona, lo que
ayuda, pero lo que retiene el calcio es el \omterm{def:b1:complexation:complex}{complejo}. Mantener el pH alto
mantiene el EDTA como \ce{Y^4-}.
\end{solution}

\begin{solution}{pb:b1:complexation:1}
\textbf{1.} \ce{Ag+ + NH3 <=> [AgNH3]+}, $\beta_1 = [\ce{AgNH3+}]/([\ce{Ag+}][\ce{NH3}])$;
\ce{Ag+ + 2NH3 <=> [Ag(NH3)2]+}, $\beta_2 = [\ce{Ag(NH3)2+}]/([\ce{Ag+}][\ce{NH3}]^2)$.
\textbf{2.} $\log K_1 = 3.32$, $\log K_2 = 7.22 - 3.32 = 3.90$.
\textbf{3.} $[\ce{AgNH3}]^{+}$ predominaría entre pNH3 3.90 y 3.32, un
intervalo invertido: no tiene dominio.
\textbf{4.} \ce{2[AgNH3]+ <=> Ag+ + [Ag(NH3)2]+}, $K = K_2/K_1 = 10^{0.58} =
3.8$.
\textbf{5.} $[\ce{Ag(NH3)2}]^{+}$ por debajo y \ce{Ag+} por encima de
$\frac12\log\beta_2 = 3.61$.
\textbf{6.} $\beta_2[\ce{NH3}]^2 = 10^{7.22 - 2.00} = 10^{5.22} = \num{1.7e5}$.
\textbf{7.} $\mathrm{p}K_d = \log\beta_2 = 7.22$.
\textbf{8.} \ce{AgCl(s) + 2NH3 <=> [Ag(NH3)2]+ + Cl-}, $K = \beta_2K_s =
10^{7.22 - 9.75} = 10^{-2.53}$.
\textbf{9.} $s = \sqrt{K}\,[\ce{NH3}] = 10^{-1.265} \times 1.0 =
\qty{0.054}{mol/L}$.
\textbf{10.} En agua, $s = 10^{-4.875} = \qty{1.3e-5}{mol/L}$: unas 4000
veces más.
\textbf{11.} $0.0543 \times 143.4 = \qty{7.8}{g}$.
\textbf{12.} $[\ce{Ag+}] = K_s/[\ce{Cl-}] = 10^{-9.75}/0.054 =
\qty{3.3e-9}{mol/L} \ll s$.
\textbf{13.} $s = \sqrt{K}(1.0 - 2s)$, $s = 0.0543/(1 + 2 \times 0.0543) =
\qty{0.049}{mol/L}$.
\textbf{14.} El amoniaco de \qty{1.0}{mol/L} tiene pH 11.6, más que
$9.25 + 1$: solo el \qty{0.4}{\percent} está como \ce{NH4+}.
\textbf{15.} $K = 10^{7.22 - 12.27} = 10^{-5.05}$.
\textbf{16.} $s = 10^{-2.525} = \qty{3.0e-3}{mol/L}$, \qty{0.56}{g/L}.
\textbf{17.} $K = 10^{-8.85}$, $s = 10^{-4.425} = \qty{3.8e-5}{mol/L}$,
\qty{8.8}{mg/L}.
\textbf{18.} Se trata cada precipitado con amoniaco diluido: solo se
disuelve el cloruro. El amoniaco concentrado disuelve el bromuro, más
despacio y en parte; el yoduro no se disuelve.
\textbf{19.} $s = 1.0/187.8 = \qty{5.3e-3}{mol/L}$; $[\ce{NH3}] = s/\sqrt{K}
= \num{5.3e-3}/10^{-2.525} = \qty{1.8}{mol/L}$.
\textbf{20.} $s = \qty{4.3e-3}{mol/L}$ y $[\ce{NH3}] = \num{4.3e-3}/10^{-4.425}
= \qty{113}{mol/L}$, el doble de la concentración del agua en el agua pura
(\qty{55}{mol/L}): imposible. El amoniaco no puede disolver el yoduro de
plata.
\textbf{21.} \ce{[Ag(NH3)2]+ + Cl- + 2H3O+ -> AgCl(s) + 2NH4+ + 2H2O},
$K = 10^{2.53} \times (10^{9.25})^2 = 10^{21.03}$: reaparece el
precipitado blanco, lo que confirma el cloruro de plata.
\textbf{22.} $s = 1.0/143.4 = \qty{6.97e-3}{mol/L}$; $[\ce{NH3}] =
s/\sqrt{K} = \num{6.97e-3}/10^{-1.265} = \qty{0.128}{mol/L}$.
\textbf{23.} $2s = \qty{0.0139}{mol/L}$.
\textbf{24.} $[\ce{Ag+}] = 10^{-9.75}/\num{6.97e-3} = \qty{2.6e-8}{mol/L}$,
despreciable frente a $s$.
\textbf{25.} $0.128 + 0.014 = $ \textbf{\qty{0.142}{mol/L} de amoniaco}:
\qty{0.142}{mol} en el litro disuelven exactamente \qty{1.0}{g} de
cloruro de plata.
\end{solution}
